\section{Stakeholders}

\colorlet{initiator}{green!40}
\colorlet{executor}{yellow!40}
\colorlet{user}{blue!40}
\colorlet{other}{orange!40}
\colorlet{affected}{red!40}

\begin{frame}{Stakeholders}
\mycomment{
    Projects, by definition, involve many persons who have not been connected in this way before.

    Everyone involved in whichever way is a stakeholder; for project management, everyone must be treated like a customer.

    Above all, it is extremely important to know just with whom you are dealing with---what these people want, what influence they have, what their needs are.

    Many projects falter because stakeholder analysis was not broad enough, and stakeholders have not been engaged enough.

    So we are talking \textbf{stakeholder analysis} and \textbf{stakeholder engagement.}
}
\mytext{
    Stakeholder is everyone who has a stake:
    
    \medskip
    
    $\rightarrow$ who has something to win (or loose)

   \medskip

   We can make out four groups:
   
    \begin{itemize}
        \item They who push it to happen: $\rightarrow$ \colorbox{initiator}{\textit{initiators}}
        \item They who do the work: $\rightarrow$ \colorbox{executor}{\textit{executors}}
        \item They who will use it: $\rightarrow$ \colorbox{user}{\textit{users}} 
        \item They who are affected in other ways: $\rightarrow$ \colorbox{other}{\textit{others}}
    \end{itemize}
}
\end{frame}

\begin{frame}{Brainstorming: Who wants it?}
    \mylonglist{
    \begin{itemize}
        \item Authorities
            \begin{itemize}
                \item different levels of Government
                \item different ministries
            \end{itemize}
        \item Business people
            \begin{itemize}
                \item those who finance the project
                \item those who realise the project
                \item those who operate the result
            \end{itemize}
        \item Users
            \begin{itemize}
                \item corporations
                \item groups
                \item individuals
            \end{itemize}
        \item Affected non-users
            \begin{itemize}
                \item may oppose it
                \item may be organised
            \end{itemize}
    \end{itemize}
    }
\end{frame}

\begin{frame}{Brainstorming: Who does it?}
    \mylonglist{
    \begin{itemize}
        \item The project manager
        \item The project management office
        \item Contractors
        \item Sub-contractors
        \item Suppliers
        \item Workers and employees
        \item The steering committee
        \item The project owners
        \item Public services
        \item Money giver? Banks?
    \end{itemize}
    }
\end{frame}

\begin{frame}{Brainstorming: Who will use it?}
    \mylonglist{
    \begin{itemize}
        \item Corporations
        \item Groups
        \item Individuals
        \item Providers, shopkeepers
        \item Customers
        \item Machines, vehicles
        \item Robots
        \item Children
        \item Squatters, homeless
        \item Animals
    \end{itemize}
    }
\end{frame}

\begin{frame}{Brainstorming: Who else are affected?}
    \mylonglist{
    \begin{itemize}
        \item Regulators, authorities
        \item Neighbours
        \item Competitors
        \item Future generations
        \item Nature
        \item Markets
        \item Gangsters
        \item Politicians
        \item The Internet
        \item ...
    \end{itemize}
    }
\end{frame}


% \begingroup
%     \setbeamercolor{background canvas}{bg=violet!20}
% \begin{frame}{Task}
%     Describe one (ideal or typical) stakeholder of your case study!
%     \begin{itemize}
%         \item Tell the story
%         \item What was / is in for them?
%     \end{itemize}
% \end{frame}
% \endgroup

\begin{frame}{Stakeholder Matrix}
    \mycomment{
    A good way to think about stakeholders is to write everyone you can think off on a card, and organise these cards in a matrix or tree.

    It is recommendable to brainstorm with everyone involved in the project at present, as stakeholders will know other stakeholders. 

    One way to sort the identified stakeholders is into the classes discussed: \colorbox{initiator}{initiators,} \colorbox{executor}{executors,} and everyone \colorbox{other}{affected}, which includes future and present \colorbox{user}{users}.
    }
    \mytable{
    \renewcommand{\arraystretch}{2}
    \begin{tabular}{|r|c|c|}
    \hline
         & internal & external \\
         \hline
         \colorbox{initiator}{initiators} & & \\
         \hline
         \colorbox{executor}{executors} & & \\
         \hline
         \colorbox{other}{affected} & & \\
         \hline
    \end{tabular}
    }
\end{frame}



\begin{frame}{Stakeholder Matrix}   
    \mycomment{
    We can arrange the stakeholders that we have identified in the step prior in a tree structure.
    }
        \showmytree[]{../shared/tikz/forest-stakeholders-suchanek-general}
\end{frame}

\begin{frame}{Stakeholder Matrix}
    \mycomment{
        Internal initiators are always a specific department, and also the top level management which backs the project.

        External initiators can be public or private persons and authorities. These initialise the project for a reason, so there are further external initiators (indirect initiators) involved, such as pressure groups, their own shareholders, the voting population, or lucrative demand.
    }
    \showmytree[ini-in,ini-ex]{../shared/tikz/forest-stakeholders-suchanek-general}
\end{frame}

\begin{frame}{Stakeholder Matrix}
    \mycomment{
    Everyone involved in making the project become reality is also a stakeholder, and imortant stakeholder. 
    
    The project management office (PMO) is core, but PMO will have to work with multiple other departments in the company, which also must be taken account of.

    Externally, subject matter experts and consultants play a significant role---or at least, they should. Contractors must be briefed and found, and their input taken into account as well. There are also always authorities directly involved in the project.
    }
        \showmytree[exe-in,exe-ex]{../shared/tikz/forest-stakeholders-suchanek-general}
\end{frame}

\begin{frame}{Stakeholder Matrix}
    \mycomment{
    The list of affected stakeholders is typically very long, and often longer than you would think at first glance. Not only is everyone who comes in sight of the project involved. The project will also cause delays while it is under way, and it will change the situation once it is completed. Therefore, again, it is necessary to think big, and not forget anyone who could be affected in the analysis, and to engage with them in one way or the other. 
    }
        \showmytree[aff-in,aff-ex]{../shared/tikz/forest-stakeholders-suchanek-general}
\end{frame}

\begin{frame}{Stakeholder Matrix}
    \mycomment{
    We can filter our lists showing only the internal stakeholders: initiators, executors, and affected. Just how are we going to engage with all these, and how in which phase of the project?

    We see that often those initialising are also affected; but not all affected are initialisers. Those affected who are not initialisers are \textbf{to be convinced.}
    }
        \showmytree[ini-in,exe-in,aff-in]{../shared/tikz/forest-stakeholders-suchanek-general}
\end{frame}

\begin{frame}{Stakeholder Matrix}
    \mycomment{
    Similarly, we can filter per external stakeholders. Again, some affected are among the initialisers, and some are not. These relations must be carefully analysed. 
    
    Those who are affected but not among the initialisers may either be negatively affected, or they may feel left out. 

    Those who are involved in all branches may then develop their own agenda, and are among the most influencial and important stakeholders.
    }
        \showmytree[ini-ex,exe-ex,aff-ex]{../shared/tikz/forest-stakeholders-suchanek-general}
\end{frame}





\begin{frame}{More on stakeholders}
\mycomment{
    Every stakeholder is different, and none of them fits exactly the stereotype. 

    Firstly, stakeholders can be directly affected, or to varying degrees. Much of that has to do with the effect of the project on those near to them: their children; their opponents; their allies. 

    Each and everyone has a bundle of interests, and may act more or less consciously. They will have very different experience both with the project and with the other stakeholders.

    Human interaction is more than the sum of individuals acting. Stakeholders engage with one another, and group dynamics can arise that may or may not be favourable for the project. 

    It is pivotal to understand properly stakeholder's interest and means to communicate.
}   
\mytext{
    \begin{itemize}
        \item directly vs. indirectly involved stakeholders
        \item owners / invested stakeholders
        \item interests of people involved
        \item roles people play
        \item \textbf{actual persons}
        \item different perspectives
        \item different expertise
        \item different power
        \item different personal position
    \end{itemize}
}
\end{frame}



\begin{frame}{Stakeholder Journey}
    \mycomment{
    As the project advances from idea to handover, stakeholders are involved in various stages.

    Each stakeholder will have an individual \textbf{stakeholder journey} from the first time they hear about the project until the end. Their stance towards the project, their involvement, their investment may change in the process.

    Clever stakeholder engagement considers these journeys and tries to create favourable stakeholder journeys.
    }
    \mytext{
    \begin{itemize}
        \item How do they get involved?
        \item First impression?
        \item First reaction?
        \item Positive side?
        \item Negative side?
        \item What is their stake?
        \item Where does the path split?
        \begin{itemize}
            \item when are they convinced?
            \item when do they change their stand?
            \item when do they `up their ante'?
            \item when do they loose interest?
        \end{itemize}
        \item What narrative do you tell them?
        \item What narrative do you want them to tell?
    \end{itemize}
    }
\end{frame}

\begin{frame}{Stakeholder Journey}
    \mycomment{
    Every stakeholder is your customer; so the concept of \textbf{customer journey} can be applied. Understand each of your stakeholders as a customer---and understand that each of them has an individual history.

    The experience should be carefully managed. How do your stakeholders know about the project? What decisions can they take---to get involved, to oppose it? How do they decide? And once they have decided, will they re-consider? Can you make them re-consider, invest more, or change their opposition? 

    It is recommendable to really write up some (imaginary or real) customer journeys, so that you have a proper understanding of where your stakeholders are coming from, where they are at, and where they are likely heading.
    }
    \showmytikz[]{../shared/tikz/tikz-customerjourney}
\end{frame}

\begin{frame}{What is in it for me?}
\mycomment{
    People usually judge projects in their own interest, that are the benefits and drawbacks the project has for themselves.

    Only to a minor extent do people take the benefits and drawbacks for other people into account. 
}
\mytext{
    People usually want to gain...
    \begin{itemize}
        \item money?
        \begin{itemize}
            \item money quick?
            \item money long term?
        \end{itemize}
        \item fame?
        \item security?
        \item idealistic win?
        \item something for those near to them?
    \end{itemize}
    
    \medskip
    Are they rational? Does that matter?
}
\end{frame}

\begin{frame}{What do I put into it?}
\mycomment{
    Stakeholders can be invested, that is that they put money into the project which they can lose.

    But stakeholers can be involved in many other ways. A good project manager should offer ways for anyone affected to input something into the project.
}
\mytext{
    \begin{itemize}
        \item money?
        \begin{itemize}
            \item when, to whom, and how
            \item to which conditions
        \end{itemize}
        \item publicity?
        \item expertise?
        \item connections?
        \item work hours 
        \begin{itemize}
            \item my own?
            \item my team members?
        \end{itemize}
    \end{itemize}
}
\end{frame}

\begin{frame}{Who can stop it?}
\mycomment{
In their stakeholder journey, some may come to actually oppose the project. 

This is to be expected, since projects deliver change, and change is not viewed positively unless it brings substantial benefits.

But not everyone opposing the projects creates serious obstacles. Thus, we must also analyse power and influence.
}
\mytext{
    \begin{itemize}
        \item Ask: who can do it?
        \item but also: who can spoil it?
        \item Questions of power and influence
    \end{itemize}
}
\end{frame}


\begin{frame}{What can I do to stop it?}
\mycomment{
    Think about not just who could oppose the project, but also by which ways. This will take centre stage again as we will be looking into risk management. 

    The better you understand what people could do to jeopardise your project, the better you can create stakeholder journeys that have people be in your favour instead.
}
\mytext{
    \begin{itemize}
        \item withdraw my support?
        \item talk bad about it / you?
        \item not do what I am supposed to do?
        \item confuse people?
        \item involve more people?
        \item go to court?
        \item ...
    \end{itemize}
}
\end{frame}

 
% \begin{frame}{Stakeholder Map / Team Map}  % depends on: forest

%     \begin{forest}
%     for tree={
%         font=\footnotesize,
%         grow = east,
%         parent anchor = east,
%         child anchor = west,
%         forked edge,
%         inner sep=2pt,
%         anchor=base west
%             }
%     [\parbox{20ex}{\centering Project: \\ Acquisition of another company}
%         [3rd parties [...][...]]
%         [Account managers[...][...]]
%         [Human resources[...][...]]
%         [The media[...][...]]
%         [Business divisions[...][...]]
%         [Marketing[...][...]]
%         [Technology[...][...]]
%         [Finance[...][...]]
%     ]
%     \end{forest}
    
% \end{frame}


\begin{frame}{Stakeholder Matrix}
    \showmytikz{../shared/tikz/tikz-stakeholdermatrix2}
    \mycomment{
    Now that we have carefully collected stakeholders with their entire profiles, we can group them by \textbf{power} and \textbf{support.} 

    If you have collected your stakeholders on cards on a (magnetic) white board as suggested earlier, you can now use these cards to sort them in this way.

    Note that instead of a matrix proper, we can be more granular and treat both power and support as axes, where clusters may appear.
    }
\end{frame}

\begin{frame}{Stakeholder Matrix}
    \showmytikz[action]{../shared/tikz/tikz-stakeholdermatrix2}
    \mycomment{
        Once we have sorted our stakeholders by power and support, we can develop strategies to deal with these.

        You should

        \begin{itemize}
            \item \textbf{keep promoters close.} Among these are your product owners and is your entire team, the board, your core clients, and so on.
            \item \textbf{keep supporters informed.} These could include top level management (if not much involved), donors, banks, parts of the public. You need not use the same channel to inform them all; we will come to this later.
            \item \textbf{keep disruptors satisfied.} Those are dangerously powerful people who are not your greatest fans. You must be careful to keep them satisfied, so that they do not begin to jeopardise the project.
            \item \textbf{keep an eye on destractors.} These are \enquote{barking dogs}---people who oppose your project, but have little power. You must not waste too much time on these, but watch out how they might influence the other groups!
        \end{itemize}
    }
\end{frame}



% \begin{frame}{No Man is an Island!}
% \mytext{
%     \begin{itemize}
%         \item Every project has multiple stakeholders
%         \item direct stake holders
%         \item owners
%         \item steering groups
%         \item project management team
%         \item team members
%         \item supervisors
%         \item consultants
%         \item potential clients
%         \item 3rd parties, indirect stakeholders
%         \item etc.
%     \end{itemize}
% }
    
% \end{frame}


% \begin{frame}[allowframebreaks]{Environment}
%     \begin{tikzpicture}
%         \draw (0,0) rectangle (5,4);
%         \node[draw] at (2.5,2) {Project};
%     \end{tikzpicture}
    
%     \framebreak
    
%     \begin{tikzpicture}
%         \draw (0,0) rectangle (5,4);
%         \draw (1,1) rectangle (4,3);
%         \node[draw] at (2.5,2) {Project};
%     \end{tikzpicture}
    
%     \framebreak
    
%     \begin{quote}
%         A successful project is best characterised by the change it delivers. \cite[p.12]{Roberts}
%     \end{quote}
% \end{frame}


% \begin{frame}{Product Flow Chart (part)}
%     \tikzstyle{decision} = [diamond, draw, fill=blue!20, 
%         text width=4.5em, text badly centered, node distance=2cm, inner sep=0pt]
%     \tikzstyle{block} = [rectangle, draw, fill=blue!20, 
%         text width=5em, text centered, rounded corners, minimum height=4em]
%     \tikzstyle{line} = [draw, -latex']
%     \tikzstyle{cloud} = [draw, ellipse,fill=red!20, node distance=3cm,
%         minimum height=2em]
        
%     \begin{tikzpicture}[scale=0.2, node distance = 2.2cm, auto]
%     \footnotesize
%         % Place nodes
%         \node [block] (case) {business case};
%         \node [block, below of=case] (gov) {governance document};
%         \node [block, below left of=gov] (use) {user requirements document};
%         \node [block, below right of=gov] (sol) {solution design document};
%         \node [block, below right of=use] (sel) {selected factory};
%         % Draw edges
%         \path [line] (case) -- (gov);
%         \path [line] (gov) -- (use);
%         \path [line] (gov) -- (sol);
%         \path [line] (use) -- (sel);
%         \path [line] (sol) -- (sel);
%     \end{tikzpicture}
    
%     Source: \cite[p. 169]{Roberts}
    
% \end{frame}




\begin{frame}{Stake holder number over time}
    \mycomment{
    As stakeholder have a live and are on a stakeholder journey, not all of them will be involved from the start, and not all of them until the end.

    However, stakeholders \enquote{accumulate}. While in the initial phases only few people know about the idea, at the end many hundreds, often thousands, if not millions, of people are involved.

    If the project takes unforseen turns, the stakeholders involved may also change radically. Think about something blowing up!
    }
    \mytext{
    \begin{itemize}
        \item Number typically grows over times
        \item `From top to bottom'
        \item Some will quit having a stake after a while (payoff)
        \item After the end of project...?
        \item `And then the lawyers get involved' 
    \end{itemize}
    }
\end{frame}

\begin{frame}{Stakeholder Involvement Strategy}
    \mytext{
    For each stakeholder group, ideally for each stakeholder, have a strategy how to communicate with them and how to engage them.

    If you don't know the person(s) already, imagine them, and write up realistic stories for realistic people.
    }
\end{frame}

\begin{frame}{Stakeholder Involvement Strategy}
\mytable{
    \scriptsize
    \begin{tabular}{>{\RaggedRight}p{18ex}>{\RaggedRight}p{33ex}>{\RaggedRight}p{20ex}}
    \toprule
     Stakeholder \& Role & Involvement & Strategy \\
     \midrule
     John Doe, Owner & 
     Came up with the project. Millionaire from the US. Has already found funding. Philanthropic motivation. & Report monthly. Provide figures.\tabularnewline
     \addlinespace
     Jane Doe, Manager & Wife of John Doe, in charge of making the project happen. Visual artist. &  Discuss PR and visuals. Give ceremonial tasks. \\
     \addlinespace
     Rick Mayer,\\ Personal Assistant to the Does & The man in charge of everything. Trained in business administration. Not qualified in content matters of the project.  & Speak to daily or as required. Provide detail on progress and requirements. \\
      ... & ... & ... \\
      ... & ... & ... \\
     \bottomrule
    \end{tabular}
}
\end{frame}


